% Supplemental proof fragment; use the same theorem environments as paper.tex.
\section*{平移质量的量化证书}
以下两个附件保留第十六、十七轮中不再作为主稿推导所需的独立几何系数。
先令 $K$ 为任意范数单位球、$\rho=\|\cdot\|_K$、
$d\mu=\rho^{\tau-n}dx$、$\kappa=n|K|/\tau$，其中 $0<\tau<n$。
\begin{lemma}\label{cert:norm-shift}
若 $C\ge2$，则
\[0\le1-\inf_{\rho(u)\le C^{-1}}\frac{\mu(u+K)}\kappa
\le\tau(n-\tau)2^{n+2-\tau}C^{-2}.\tag{TC.1}\]
\end{lemma}
\begin{proof}
令 $\delta=C^{-1}$。层析表示为
$\mu(A)=(n-\tau)\int_0^\infty|A\cap sK|s^{\tau-n-1}\,ds$。
中心与平移球的交叠体积差非负，仅在 $s\in[1-\delta,1+\delta]$ 内非零。
壳内两交集都包含 $(1-\delta)K$，故差至多 $n|K|\delta$。
利用 $s^{\tau-n-1}\le2^{n+1-\tau}$ 及壳宽 $2\delta$，再除以
$\kappa$ 即得界。该量化界与范数球形状无关，但不声称常数最优。
\end{proof}
下面专取 $Q(x,r)=x+[-r,r]^n$、$d\mu=\|x\|_\infty^{\tau-n}dx$，
$\kappa=n2^n/\tau$，并令
\[A_C=\inf_{\|u\|_\infty\le C^{-1}}\mu(Q(u,1))/\kappa.\]
\begin{lemma}[平移损失]\label{cert:cube-shift}
当 $C\to\infty$ 时，
\[A_C=1-\frac{\tau(n-\tau)}2C^{-2}+O_{n,\tau}(C^{-3}).\tag{TC.2}\]
\end{lemma}
\begin{proof}
层析恒等式给出
\[\mu(A)=(n-\tau)\int_0^\infty |A\cap Q(0,s)|s^{\tau-n-1}\,ds.\tag{TC.3}\]
令 $\delta=C^{-1}$。每个坐标的区间交叠长度随平移的绝对值单调不增，
故 $A_C$ 在 $u=(\delta,\ldots,\delta)$ 取得。相应交叠方体的边长为
\[L_\delta(s)=\begin{cases}
2s,&0\le s\le1-\delta,\\
s+1-\delta,&1-\delta\le s\le1+\delta,\\
2,&s\ge1+\delta.
\end{cases}\]
中心与平移方体的交叠体积差只在上述薄壳内非零。令 $s=1+\delta v$，
在 $-1\le v\le1$ 上一致地有
\[(2\min\{1,s\})^n-L_\delta(s)^n
=n2^{n-1}\delta(1-|v|)+O_n(\delta^2).\]
因壳远离零，$s^{\tau-n-1}=1+O_{n,\tau}(\delta)$，而
$\int_{-1}^1(1-|v|)\,dv=1$。代入 (TC.3) 得
\[\mu(Q(0,1))-\mu(Q(u,1))
=(n-\tau)n2^{n-1}\delta^2+O_{n,\tau}(\delta^3).
\]
除以 $\kappa=n2^n/\tau$ 即得 (TC.2)。这里未在权的内部奇点处求导。
\end{proof}
